\section{Preserving six permanent copies under internal padding}

\subsection*{Finding F-026: a fixed-entry coefficient projection and an all-copy injection}
Let \(n\geq3\), and let \(\nu\vdash3d\) satisfy
\(\ell(\nu)\leq7\) and \(\nu_3\leq2\). Define
\[
                    \nu[n]=(\nu_1+d(n-3),\nu_2,\ldots).
\]
For the internally padded \(3\)-by-\(3\) permanent,
\begin{equation}\label{eq:F026copies}
                   q^3_{\nu[n]}(d[n])\geq a_\nu(d[3]).
\end{equation}
In particular, for every \(n\geq3\),
\begin{equation}\label{eq:F026six}
              q^3_{(12n-17,7,2^5)}(12[n])\geq6.
\end{equation}
The right side is a count of independent actual permanent
copies. An equality with the full ambient multiplicity
at inner degree \(n\) is not asserted when \(\nu_2>3\).
This distinction applies in particular to \eqref{eq:F026six}.

\paragraph{The coefficient map with its distinguished entry fixed.}
Use the same seven variables, target \(\Gamma\),
coefficient ordering and source numbering as in F-025.
Require the first matrix entry to equal \(x_1\), so
only the 56 source coordinates numbered 7 through 62
may vary. The following new integer pencil meets that
condition:
\begin{center}
\begin{tabular}{@{}c|rrrrrrr@{}}
\toprule
\((i,j)\)&1&2&3&4&5&6&7\\\midrule
\((1,1)\)&1&0&0&0&0&0&0\\
\((1,2)\)&0&4&2&-3&-1&0&-3\\
\((1,3)\)&2&1&3&4&1&3&-3\\
\((2,1)\)&0&3&4&-1&2&-2&2\\
\((2,2)\)&-3&-4&-2&-1&1&-2&3\\
\((2,3)\)&3&3&0&2&-3&4&2\\
\((3,1)\)&-2&-3&-3&3&1&-1&3\\
\((3,2)\)&0&-1&2&2&1&4&-1\\
\((3,3)\)&4&-3&-1&-3&-1&4&-1\\\bottomrule
\end{tabular}
\end{center}
The Jacobian is again specified by the coefficient formula
(F-025.4), now at this pencil. Retain the 49 allowed columns
\[
                 C'=\{7,8,\ldots,47\}
                        \cup\{49,50,\ldots,54\}\cup\{56,57\}.
                                                               \tag{F-026.1}
\]
Their square determinant is
\[
 \begin{gathered}
       \Delta'=31568108869448599905478876197553019094163200564336768,
          \\
          \Delta'\equiv641253\pmod{1000003}.
 \end{gathered}                                      \tag{F-026.2}
\]
Thus the projection onto all 49 coordinates in \(\Gamma\)
is dominant even with \(A_{11}=x_1\) fixed.
The algebraic argument is exactly the invertible-differential
argument in F-025: restrict to these 49 directions, translate,
and normalize the linear term to obtain \(t+O(t^2)\).
It works over \(\mathbb Z[(\Delta')^{-1}]\).

The complete finite witness is the displayed table together
with \(\Gamma\), (F-025.4), and (F-026.1).
The scripts
\path{agents/jet_review/check_per3_fixed_entry_transverse.py}
and
\path{agents/jet_review/verify_per3_fixed_entry_transverse.py}
construct the cofactor matrix and independently expand the
full permanent. They verify all coefficient columns and
compute (F-026.2) exactly. The second calculation includes
the seven columns which are subsequently held fixed, but
the selected minor uses only the 56 permitted columns.

\paragraph{The highest-weight injection and the actual padding pullback.}
Let \(U\) be the coordinate-function space on the first
seven variables. For a nonzero ambient type, let \(H_\nu\)
be its entire highest-weight multiplicity space in
\(\operatorname{Sym}^d\operatorname{Sym}^3 U\).
Multiplication in each column of the Cauchy matrix model
by \(e_1^{n-3}\) gives an injective map
\[
 T_n:H_\nu\longrightarrow
           \bigl(\operatorname{Sym}^d\operatorname{Sym}^n U\bigr)
                          ^{\mathrm{HW},\nu[n]}.
                                                               \tag{F-026.3}
\]
Indeed the upper unipotent group fixes \(e_1\), the
weight changes by \(d(n-3)e_1\), and column permutations
commute with this multiplication. The map is injective on
monomial tensors and therefore on their invariant subspace.
Only an injection is used here; F-014's extra conditions
for surjectivity are not assumed.

Let \(R_n\) be the actual pullback induced by padding a
cubic form \(F\) as \(x_1^{n-3}F\).
Under the natural differential, or apolar, pairing, on a
cubic monomial generator \(e^\alpha\) one has
\[
       (R_nT_n)(e^\alpha)
          =\frac{(\alpha_1+n-3)!}{\alpha_1!}\,e^\alpha.
                                                               \tag{F-026.4}
\]
The factors in the other coordinates cancel. For products
of \(d\) generators this formula extends multiplicatively.
Thus \(R_nT_n\) is an invertible diagonal operator on
the relevant polynomial coefficient algebra over
\(\mathbb C\). This is the same precise padding
calculation as O-012 and F-015; it is not an identification
of the highest-weight lift with the padding morphism itself.

\begin{proof}[Every original ambient copy remains independent]
By F-016, every polynomial in \(H_\nu\) uses only cubic
coefficient generators whose suffix degree from the third
variable is at most \(\nu_3\leq2\). Hence
\(H_\nu\subset\mathbb C[\Gamma]\). The diagonal
operator (F-026.4) preserves this coefficient algebra and
is invertible on it. It need not preserve its highest-weight
subspace; that is unnecessary.

Suppose a linear combination of the highest-weight vectors
\(T_n(H_\nu)\) vanishes on the internally padded permanent
closure. Evaluate it on the substitutions furnished by the
fixed-entry pencils above. Because \(A_{11}=x_1\), these
are exactly
\[
                         x_1^{n-3}\operatorname{per}_3 A(x),
\]
with the padding form equal to the substituted distinguished
matrix entry. They are points of the required orbit closure
by endomorphism substitution in its \(n^2\)-variable space.
Pulling the alleged linear relation back gives the
corresponding combination of \((R_nT_n)(H_\nu)\)
vanishing on the dominant \(\Gamma\)-projection of the
fixed-entry permanent family. Injectivity of that projection's
coordinate pullback makes this combination zero as a polynomial
in \(\mathbb C[\Gamma]\). The invertibility of (F-026.4)
then makes the original combination in \(H_\nu\) zero.

Therefore the \(a_\nu(d[3])\) highest-weight vectors in
the image of (F-026.3) remain independent on the padded
permanent closure. This proves \eqref{eq:F026copies}.
For \(\nu=(19,7,2^5)\), the known and independently
checked ambient coefficient is six by F-025. Its first
part shifts by \(12(n-3)\), giving \eqref{eq:F026six}.
\end{proof}

If additionally \(\nu_2\leq3\), the all-copy stabilization
F-014 identifies the new ambient multiplicity with the old
one, so \eqref{eq:F026copies} is an equality with both.
For the second row seven in \eqref{eq:F026six}, no such
conclusion follows from this argument.

\paragraph{The shifted obstruction stops at inner degree three.}
For the particular family \(\nu[n]=(12n-17,7,2^5)\),
O-027 gives full determinant multiplicity for every \(n\geq4\):
its length is seven, \(7\leq2n\), and its third row is two.
Consequently
\[
       6\leq q^3_{\nu[n]}(12[n])
                 \leq K_n(\nu[n])=a_{\nu[n]}(12[n])
                           \qquad(n\geq4).
                                                               \tag{F-026.5}
\]
Thus the exact \(6>5\) obstruction at \(n=3\) from F-025
does not persist under these first-row shifts. Six actual
permanent copies do persist, and the determinant closure
has enough copies to accommodate them at every larger
inner degree. The finite independent determinant-Jacobian
certificates in this folder also verify the seven-variable
projection at \(n=4,5\), with ranks 70 and 91 respectively;
the complete arbitrary-size proof of (F-026.5) is O-027.

\paragraph{Grothendieck dictionary and scope.}
The fixed-entry pencil parameter space is an affine linear
subscheme defined over \(\mathbb Z\), and the coefficient
projection is dominant after base change to \(\mathbb C\),
with an integral coordinate injection after inverting
\(\Delta'\). The highest-weight map is multiplication
by the first basis line in each column; the Schur-functor
and Cauchy interpretation is SGA~6, VI, Theorem~3.3,
pp.~392--393, together with Macdonald I(4.3), p.~63,
Appendix~A(5.3)--(5.6), pp.~155--156.
The precise torus grading and ideal-comodule statements
are SGA~3, I, Propositions~4.7.2--4.7.3, pp.~34--36.
The factorial comparison is used in characteristic zero;
it is not silently asserted to be invertible in every
residue characteristic. The theorem counts actual surviving
copies and does not infer an actual obstruction from their
positivity alone.
